% Phụ lục U — Bộ chính sách công khai cho người dùng. Bảng tổng hợp ở đầu dẫn
% căn cứ pháp lý; toàn văn ba chính sách viết để đọc độc lập như Phụ lục D.

\section*{Phụ lục U. Bộ chính sách}
\label{apdx:chinhsach}
\addcontentsline{toc}{section}{Phụ lục U. Bộ chính sách}

Điều khoản dịch vụ ở Phụ lục D là quy chế chung. Bên cạnh đó, Luật Thương mại
điện tử 2025 (Điều 11) và Nghị định 248/2026/NĐ-CP (Điều 4--16) yêu cầu nền
tảng có chức năng đặt hàng công khai một bộ chính sách với nội dung tối thiểu
cho từng loại \parencite{luattmdt2025, nd248_2026}. Bảng~\ref{tab:bochinhsach}
liệt kê mười chính sách Farmery công bố, đúng bằng danh sách trang con của mục
``Chính sách'' trên ứng dụng mua hàng. Ba chính sách tác động trực tiếp nhất
tới trục bán lẻ --- kiểm hàng, đổi trả và hoàn tiền, bảo mật dữ liệu cá nhân
--- được trình bày toàn văn ngay sau bảng, viết để người dùng đọc độc lập.

\begingroup
\tblsetup\tblzebra
\begin{longtable}{L{3.30cm} L{2.70cm} L{6.20cm} L{2.70cm}}
\caption{Bộ chính sách công khai của Farmery}
\label{tab:bochinhsach} \\
\toprule
\rowcolor{tblheadbg}
\textbf{Chính sách} & \textbf{Áp dụng cho} & \textbf{Nội dung chính} & \textbf{Căn cứ} \\
\midrule
\endfirsthead
\toprule
\rowcolor{tblheadbg}
\textbf{Chính sách} & \textbf{Áp dụng cho} & \textbf{Nội dung chính} & \textbf{Căn cứ} \\
\midrule
\endhead
\bottomrule
\endfoot
Quy chế hoạt động (Điều khoản dịch vụ) & Mọi người dùng & Thông tin chủ quản, quyền và nghĩa vụ các bên, phân định trách nhiệm theo hình thức bán, xử lý vi phạm (Phụ lục D) & Luật TMĐT Đ11; NĐ 248 Đ4, Đ6 \\
Kiểm hàng khi nhận & Khách lẻ & Được kiểm gì, cách từ chối nhận tại cửa, ghi nhận minh chứng; toàn văn ở mục U.1 & NĐ 248 Đ13 khoản 5 \\
Đổi trả và hoàn tiền & Khách lẻ; khách sỉ theo hợp đồng & Điều kiện, thời hạn theo nhóm hàng, quy trình, phương thức hoàn, chi phí hoàn trả; toàn văn ở mục U.2 & Luật TMĐT Đ11; NĐ 248 Đ14 \\
Hủy đơn & Khách lẻ & Hủy miễn phí trước khi đóng gói; sau đó xử lý theo đổi trả; hoàn tiền nếu đã trả trước & NĐ 248 Đ10 \\
Giao hàng & Khách lẻ & Giờ chốt đơn, chuyến xe lạnh về TP.HCM, khung giờ giao nội thành, vùng giao, thời gian ước tính, phí giao và ngưỡng miễn phí, tên đơn vị vận chuyển & NĐ 248 Đ9, Đ13 \\
Thanh toán & Khách lẻ, khách sỉ & Phương thức thanh toán, điều kiện trả tiền khi nhận hàng, cách hoàn tiền, đơn vị trung gian thanh toán & NĐ 248 Đ10 \\
Bảo mật dữ liệu cá nhân & Mọi người dùng & Mục đích, phạm vi thu thập và sử dụng, thời hạn lưu, ai được tiếp cận, biện pháp bảo vệ, quyền xem, sửa, xóa; toàn văn ở mục U.3 & Luật BVDLCN 2025; NĐ 248 Đ5 \\
Khiếu nại và giải quyết tranh chấp & Mọi người dùng & Kênh tiếp nhận trực tuyến gắn với đơn hàng, các bước xử lý, thời hạn phản hồi ban đầu 24 giờ và thời hạn quyết định 3 ngày làm việc, hòa giải & NĐ 248 Đ7 \\
Nghiệm thu và thanh toán thu mua & Nhà cung cấp bán cho Farmery & Tiêu chuẩn nghiệm thu theo nhóm hàng, xử lý hàng không đạt, giá điều chỉnh, thời điểm thanh toán sau nghiệm thu & Thỏa thuận thu mua (Phụ lục D Điều 8) \\
Thu hồi sản phẩm & Mọi người dùng & Khi nào thu hồi, cách thông báo người đã mua, cảnh báo trên trang truy xuất, đổi trả và hoàn tiền & Luật TMĐT Đ16, Đ17; NĐ 37/2026 Đ52 \\
\end{longtable}
\endgroup

Ngoài mười chính sách trên, chính sách giá và biểu phí cùng tiêu chí ưu tiên
hiển thị cũng được công khai theo yêu cầu của Nghị định 248/2026 (Điều 8 và
Điều 11); nội dung của chúng đã nêu tại Điều 7 và Điều 9 của Phụ lục D.

\subsubsection*{U.1. Chính sách kiểm hàng khi nhận}

\noindent\textbf{1. Phạm vi.} Chính sách áp dụng cho mọi đơn hàng do Farmery
bán và giao tới người mua.

\noindent\textbf{2. Người mua được kiểm gì.} Khi nhân viên giao hàng tới, người
mua được mở túi ngoài để xem tình trạng hàng: số lượng gói, tem truy xuất trên
từng gói, hàng có bị dập, úng, héo hay chảy nước hay không, và nhiệt độ cảm
nhận với hàng cần giữ lạnh. Người mua không bóc tem niêm phong của từng gói
trước khi nhận, trừ khi gói đã rách hoặc hàng hỏng nhìn thấy được qua bao bì.

\noindent\textbf{3. Từ chối nhận tại cửa.} Nếu hàng hỏng, thiếu hoặc không đúng
đơn, người mua được từ chối nhận toàn bộ hoặc một phần. Nhân viên giao hàng
chụp ảnh tình trạng hàng và ghi lý do trên ứng dụng giao nhận. Farmery hoàn
toàn bộ tiền phần hàng bị từ chối và cước giao, theo thời hạn của chính sách
đổi trả và hoàn tiền; người mua không phải gửi thêm minh chứng.

\noindent\textbf{4. Khi không kiểm được tại cửa.} Nếu người mua nhờ người khác
nhận hoặc yêu cầu để hàng tại nơi nhận, việc kiểm hàng được thay bằng quyền
khiếu nại trong thời hạn của chính sách đổi trả và hoàn tiền, kèm ảnh hoặc
video tình trạng hàng.

\noindent\textbf{5. Từ chối không có lý do về chất lượng.} Người mua từ chối
nhận đơn trả tiền khi nhận hàng mà không có lý do về chất lượng thì tài khoản
không còn được dùng hình thức trả tiền khi nhận hàng cho các đơn sau.

\subsubsection*{U.2. Chính sách đổi trả và hoàn tiền}

\noindent\textbf{1. Phạm vi.} Chính sách áp dụng cho hàng do Farmery bán cho
người tiêu dùng. Đơn hàng số lượng lớn do nhà cung cấp bán được đổi trả theo
hợp đồng giữa hai bên; Farmery hỗ trợ hòa giải khi có yêu cầu.

\noindent\textbf{2. Điều kiện được đổi trả, hoàn tiền.}
\begin{itemize}
    \item Hàng hỏng, dập, úng, héo vượt mức bình thường của nông sản tươi.
    \item Hàng không đúng loại, không đúng số lượng hoặc khối lượng đã đặt.
    \item Hàng thuộc lô bị thu hồi.
    \item Hàng khô còn nguyên niêm mà người mua đổi ý.
\end{itemize}
Không áp dụng đổi ý đối với hàng tươi, vì hàng đã rời chuỗi bảo quản thì không
thể bán lại cho người khác.

\noindent\textbf{3. Thời hạn gửi yêu cầu}, tính từ lúc đơn được giao thành công:
\begin{itemize}
    \item Rau ăn lá, rau gia vị, nấm tươi: 24 giờ.
    \item Củ, quả tươi: 48 giờ.
    \item Hàng khô (nông sản sấy, đóng gói kín): 7 ngày.
\end{itemize}
Các mốc thời hạn có thể được điều chỉnh; thay đổi được công bố trước và không
áp dụng ngược cho đơn đã đặt.

\noindent\textbf{4. Quy trình.}
\begin{enumerate}
    \item Người mua gửi yêu cầu ngay trên trang đơn hàng, chọn sản phẩm và lý
    do, kèm ảnh hoặc video.
    \item Farmery xác nhận đã tiếp nhận trong vòng 24 giờ. Mọi trao đổi tiếp
    theo diễn ra trong luồng trao đổi của chính yêu cầu đó.
    \item Farmery ra quyết định trong vòng 3 ngày làm việc kể từ khi tiếp nhận:
    hoàn tiền toàn phần, hoàn một phần, giao bù sản phẩm tương đương, hoặc từ
    chối khi bằng chứng không đủ; quyết định luôn kèm lý do.
    \item Người mua không đồng ý với quyết định có thể yêu cầu xem xét lại một
    lần, hoặc chuyển sang giải quyết tranh chấp theo Điều 16 của Điều khoản
    dịch vụ.
\end{enumerate}

\noindent\textbf{5. Phương thức hoàn tiền.} Tiền được hoàn qua đúng phương
thức người mua đã thanh toán; đơn trả tiền khi nhận hàng được hoàn qua tài
khoản ngân hàng người mua cung cấp. Thời hạn hoàn tiền là 7 ngày làm việc kể
từ khi có quyết định.

\noindent\textbf{6. Trả hàng và chi phí.}
\begin{itemize}
    \item Khi giá trị yêu cầu dưới mức Farmery công bố, người mua được hoàn
    tiền mà không phải gửi trả hàng.
    \item Khi phải trả hàng vì lỗi của Farmery hoặc của đơn vị vận chuyển,
    Farmery chịu toàn bộ chi phí thu hồi.
    \item Khi trả hàng khô do đổi ý, người mua chịu cước gửi trả.
\end{itemize}

\noindent\textbf{7. Lỗi của đơn vị vận chuyển.} Farmery giải quyết với người
mua trước theo chính sách này, rồi tự làm việc với đơn vị vận chuyển; người
mua không phải liên hệ đơn vị vận chuyển.

\subsubsection*{U.3. Chính sách bảo mật dữ liệu cá nhân}

\noindent\textbf{1. Dữ liệu được thu thập và mục đích.}
\begin{itemize}
    \item Thông tin tài khoản (họ tên, số điện thoại, email) và sổ địa chỉ:
    để đăng nhập, liên lạc và giao hàng.
    \item Lịch sử đơn hàng, thanh toán, khiếu nại và trao đổi: để thực hiện
    giao dịch, xuất hóa đơn, xử lý khiếu nại.
    \item Với nhà cung cấp: giấy tờ định danh, ảnh chân dung khi xác thực, mã
    số thuế, loại người bán, thông tin vùng trồng, tài khoản nhận tiền: để xác
    thực người bán theo quy định, thanh toán và thực hiện nghĩa vụ thuế.
    \item Với doanh nghiệp mua hàng: giấy phép kinh doanh, mã số thuế, người
    đại diện: để xác minh doanh nghiệp và lập hợp đồng.
    \item Từ khóa tìm kiếm và thao tác trên ứng dụng: để cải thiện kết quả tìm
    kiếm và gợi ý, ở dạng ẩn danh hoặc tổng hợp khi không cần gắn với người
    dùng cụ thể.
\end{itemize}
Farmery chỉ thu thập dữ liệu cần cho các mục đích trên, không bán dữ liệu cá
nhân cho bên thứ ba.

\noindent\textbf{2. Ai được tiếp cận.} Nhân sự Farmery chỉ được xem dữ liệu cần
cho đúng công việc được giao; mọi thao tác trên dữ liệu người dùng đều được
ghi nhật ký. Dữ liệu được chia sẻ cho đối tác ở mức tối thiểu cần thiết: tên,
số điện thoại, địa chỉ cho đơn vị vận chuyển; thông tin thanh toán cho đơn vị
trung gian thanh toán; thông tin hóa đơn cho tổ chức cung cấp dịch vụ hóa
đơn; và cho cơ quan nhà nước khi có yêu cầu theo quy định của pháp luật.

\noindent\textbf{3. Thời hạn lưu.} Dữ liệu tài khoản được lưu trong thời gian
tài khoản còn hoạt động. Dữ liệu hợp đồng được lưu ít nhất 3 năm kể từ khi
giao kết; hóa đơn, chứng từ và dữ liệu nhãn hàng hóa kèm dữ liệu giao dịch
được lưu ít nhất 5 năm. Ảnh giấy tờ định danh dùng để xác thực được lưu tới
khi tài khoản nhà cung cấp chấm dứt và hết thời hạn giải quyết khiếu nại liên
quan, sau đó bị xóa.

\noindent\textbf{4. Biện pháp bảo vệ.} Mật khẩu được băm một chiều; thông tin
định danh, mã số thuế, tài khoản ngân hàng và các tệp giấy tờ được mã hóa khi
lưu; ảnh giấy tờ đã được xóa thông tin vị trí chụp; tệp riêng tư chỉ mở được
qua đường dẫn tạm thời có thời hạn ngắn. Farmery không lưu số thẻ thanh toán.
Tài khoản nhân sự Farmery bắt buộc xác thực hai lớp.

\noindent\textbf{5. Quyền của người dùng và cách thực hiện.} Trong mục ``Tài
khoản'' của ứng dụng, người dùng có thể:
\begin{itemize}
    \item xem và tải về dữ liệu cá nhân của mình;
    \item sửa thông tin không chính xác;
    \item rút lại sự đồng ý đối với các mục đích không bắt buộc;
    \item yêu cầu xóa dữ liệu hoặc hạn chế xử lý.
\end{itemize}
Yêu cầu xóa được thực hiện trong thời hạn pháp luật về bảo vệ dữ liệu cá nhân quy định; phần dữ liệu pháp luật bắt buộc phải
lưu (mục 3) được ẩn danh thay vì xóa, và người dùng được thông báo rõ phần nào
còn được giữ lại, vì sao.

\noindent\textbf{6. Khiếu nại về dữ liệu cá nhân.} Người dùng gửi khiếu nại
qua mục ``Tài khoản'' hoặc kênh hỗ trợ đã công bố. Farmery xác nhận tiếp nhận
trong 24 giờ và trả lời trong 3 ngày làm việc. Khi xảy ra sự cố lộ lọt dữ
liệu, Farmery thông báo cho người dùng bị ảnh hưởng và cơ quan có thẩm quyền
theo quy định.

\noindent\textbf{7. Thay đổi chính sách.} Mỗi phiên bản chính sách có ngày
hiệu lực. Khi thay đổi, Farmery thông báo trước và yêu cầu người dùng đồng ý
lại nếu mục đích sử dụng dữ liệu được mở rộng.
